% problem-000200 3 海水锂选择性提取
\section*{3. 海水锂选择性提取}
随着锂电池的⼴泛应⽤，锂已成为重要的战略资源，被称为⽩⾊⽯油。据估计，海⽔中锂的总含量为陆地锂总含量的5000倍以上，但海⽔中锂的质量浓度仅为0.1\~0.2 ppm，从海⽔提取锂⾸先需要对低浓度的Li+进⾏选择性富集，Li+能够嵌⼊某些氧化物并在⼀定条件下脱出，据此可以进⾏Li+的富集。

\subsection*{3-1}
以 $\mathrm { L i _ { 2 } C O _ { 3 } }$ 和 $\mathrm { M n O _ { 2 } }$ 为原料，充分混合后在 $7 2 0 ~ ^ { \circ } \mathrm { C }$ 下煅烧3 h，冷⾄室温后即可得到复合氧化物$\mathrm { L i M n _ { 2 } O _ { 4 } }$ (LMO)。⽤ 1 mol L−1的 HCl 在 $6 0 ~ ^ { \circ } \mathrm { C }$ 处理LMO，将其中所有Li+置换后得到HMO，HMO可以和Li+反应再⽣成LMO，然后LMO与酸作⽤脱出Li+从⽽实现Li+富集。如此，可以循环处理。

\subsection*{3-1}
-1 写出合成 $\mathrm { L i M n _ { 2 } O _ { 4 } }$ 的⽅程式。

\subsection*{3-1}
-2 ⽤酸处理LMO时，除离⼦交换反应之外也会发⽣⼀个副反应，该副反应导致固体中Mn的平均氧化态有所升⾼，写出副反应对应的化学⽅程式。这⼀副反应对再⽣后的HMO的富集性能有何影响？

\subsection*{3-2}
利⽤电化学富集也是有效的⽅法之⼀。某电化学系统如图所⽰，其包含两个电池单元，中间区域置有具有⼀维孔道结构的λ- $\cdot \mathrm { M n O _ { 2 } }$ ，该孔道可以容纳合适的阳离⼦进出 $\mathrm { M n O } _ { 2 \circ }$ 电极通过阳离⼦交换膜与两个电化学池中的对电极隔开，腔室1和2中的电解质分别为0.5M的$\mathrm { N a } _ { 2 } \mathrm { S O } _ { 4 }$ 和 0.1 M 的 $\mathrm { L i O H _ { \circ } }$ 。该电化学系统的⼯作步骤如下：1）向 $\mathrm { M n O _ { 2 } }$ 所在腔室通⼊海⽔，启动电源1，使海⽔中的Li+进⼊ $\mathrm { M n O _ { 2 } }$ 结构⽽形成 $\mathrm { L i } _ { x } \mathrm { M n } _ { 2 } \mathrm { O } _ { 4 } ;$ 2）关闭电源1和海⽔通道，启动电源2，同时向电极2上通⼊空⽓，使 $\mathrm { L i } _ { x } \mathrm { M n } _ { 2 } \mathrm { O } _ { 4 }$ 中的Li+脱出进⼊腔室2。

\subsection*{3-2}
-1 为衡量 $\mathrm { M n O _ { 2 } }$ 对 Li+的富集效果，将 0.50M 的 LiCl溶液通⼊MnO_{2} (4.8 mg)所在腔室，启动电源1，使电流

（图：）

恒定在5.0mA，累计⼯作325s后发现 $\mathrm { M n O _ { 2 } }$ 的电极电势快速下降，计算所得 $\mathrm { L i } _ { x } \mathrm { M n } _ { 2 } \mathrm { O } _ { 4 }$ 中的�。3-2-2 写出上述过程中发⽣腔室2中阴极和阳极上的电极反应。
